\chapter{Ruang Vektor dan Subruang}

Tabel berikut memuat tiga rekaman aliran jaringan; setiap baris satu rekaman, setiap kolom satu
fitur (paket masuk, paket keluar, durasi, koneksi gagal). Kolom-kolomnya, dibaca dari atas ke
bawah, adalah
\[
  \mathbf{c}_{1}=(1,2,3),\quad \mathbf{c}_{2}=(2,4,6),\quad
  \mathbf{c}_{3}=(0,1,2),\quad \mathbf{c}_{4}=(1,3,5),
\]
sedangkan baris-barisnya $\mathbf{r}_{1}=(1,2,0,1)$, $\mathbf{r}_{2}=(2,4,1,3)$, dan
$\mathbf{r}_{3}=(3,6,2,5)$.

Berapa kolom yang benar-benar membawa informasi baru? Menghitung kolom tak nol memberi empat,
dan hitungan itu keliru: kolom kedua dua kali kolom pertama, kolom keempat jumlah kolom
pertama dan ketiga. Tabel itu hanya memuat dua arah data mandiri. Kalimat ``dapat dihitung
ulang'' baru punya arti setelah kita dapat mengukur berapa vektor minimal yang membangunnya.

Bab 1 sampai Bab 5 menyiapkan bahasa himpunan, aksioma bilangan real, lalu limit barisan dan
deret; semuanya bekerja dengan kedekatan. Persoalan di atas tidak memuat kedekatan sama sekali:
entrinya eksak, dan pertanyaannya murni aljabar. Himpunan mana yang tertutup terhadap kedua
operasi, dan berapa vektor paling sedikit yang membangunnya. Satu kebiasaan dari Bab 5 tetap
terpakai: di sini setiap kombinasi berhingga banyaknya.

Ada pertanyaan kedua yang menuntut alat yang sama. Penilai linear memberi skor
$s(\mathbf{x})=\mathbf{w}\cdot\mathbf{x}$ dengan $\mathbf{w}=(1,1,1)$. Arah
$\boldsymbol{\delta}$ yang tidak mengubah skor memenuhi
$\mathbf{w}\cdot\boldsymbol{\delta}=0$, yaitu sebuah bidang; setiap perturbasi terpecah
menjadi bagian yang mengubah skor dan bagian yang tidak. Apakah pecahan itu tunggal? Ya, dan
sifat itu bernama jumlah langsung.

\begin{untukanak}
Bayangkan sekotak cat poster berisi empat kaleng: merah, kuning, biru, dan jingga. Kita
boleh mencampur isi kaleng dengan takaran apa pun, dan hasil campurannya boleh dipakai
untuk mencampur lagi. Jingga bisa dibuat dari merah dan kuning. Jadi kaleng jingga tidak
menambah satu kemampuan pun yang belum ada di kotak itu; ia hanya memakan tempat.

Pertanyaan yang membuat penasaran: berapa kaleng paling sedikit yang harus tetap ada agar
semua warna yang bisa kita buat tetap bisa kita buat? Jawabannya tidak sama dengan banyaknya
kaleng yang tersisa setelah kita membuang satu, dan urutan membuangnya tidak boleh mengubah
angka itu. Itu bukan kebetulan, dan membuktikannya tidak sesederhana kelihatannya.

Ada pertanyaan yang lebih aneh. Misalkan dua orang mencampur dengan resep berbeda, lalu
hasilnya sama persis. Apakah resepnya sama? Kadang ya, kadang tidak. Di bab ini kita akan tahu
mengapa angka itu punya nama sendiri, mengapa ia tidak bergantung pada urutan membuang, dan
mengapa dua campuran berbeda kadang menghasilkan warna yang sama.
\end{untukanak}

\section{Aksioma ruang vektor dan contohnya}
\tingkat{dasar}

Ruang vektor adalah himpunan yang anggotanya boleh dijumlahkan dan boleh diperbesar atau
diperkecil dengan bilangan, dengan hasil yang tetap berada di dalam himpunan itu, dan dengan
aturan hitung yang sama seperti pada pasangan bilangan biasa.

\begin{definisi}[Ruang vektor]
\index{ruang vektor}
Himpunan tak kosong $V$ bersama operasi penjumlahan $V\times V\to V$ dan perkalian skalar
$\mathbb{R}\times V\to V$ disebut \emph{ruang vektor atas $\mathbb{R}$} jika untuk semua
$\mathbf{u},\mathbf{v},\mathbf{w}\in V$ dan semua $a,b\in\mathbb{R}$ berlaku:
\begin{enumerate}
  \item $\mathbf{u}+\mathbf{v}=\mathbf{v}+\mathbf{u}$;
  \item $(\mathbf{u}+\mathbf{v})+\mathbf{w}=\mathbf{u}+(\mathbf{v}+\mathbf{w})$;
  \item ada $\mathbf{0}\in V$ dengan $\mathbf{v}+\mathbf{0}=\mathbf{v}$ untuk setiap
        $\mathbf{v}$;
  \item ada $-\mathbf{v}\in V$ dengan $\mathbf{v}+(-\mathbf{v})=\mathbf{0}$;
  \item $a(\mathbf{u}+\mathbf{v})=a\mathbf{u}+a\mathbf{v}$;
  \item $(a+b)\mathbf{v}=a\mathbf{v}+b\mathbf{v}$;
  \item $a(b\mathbf{v})=(ab)\mathbf{v}$;
  \item $1\mathbf{v}=\mathbf{v}$.
\end{enumerate}
\end{definisi}

Kedua operasi itu membawa syarat ketertutupan di dalam tipenya: hasilnya selalu di $V$, jadi
definisi ini tidak memuat aksioma penutupan tersendiri. Delapan aksioma itu juga tidak saling
bebas, dan kita tidak mengejar daftar terpendek; yang dibutuhkan akibat-akibat kecilnya,
semuanya diturunkan tanpa melihat bentuk anggota $V$.

\begin{lemma}
Dalam sembarang ruang vektor, vektor nol tunggal, invers penjumlahan setiap vektor tunggal,
$0\mathbf{v}=\mathbf{0}$, $(-1)\mathbf{v}=-\mathbf{v}$, dan $c\mathbf{v}=\mathbf{0}$
mengakibatkan $c=0$ atau $\mathbf{v}=\mathbf{0}$.
\end{lemma}

\par\medskip\noindent\textbf{Bukti.}\ 
Untuk ketunggalan nol, misalkan $\mathbf{0}$ dan $\mathbf{0}'$ keduanya netral. Maka
$\mathbf{0}=\mathbf{0}+\mathbf{0}'=\mathbf{0}'$. Untuk ketunggalan invers, misalkan
$\mathbf{v}+\mathbf{a}=\mathbf{0}$ dan $\mathbf{v}+\mathbf{b}=\mathbf{0}$. Maka
\[
  \mathbf{a}=\mathbf{a}+\mathbf{0}=\mathbf{a}+(\mathbf{v}+\mathbf{b})
  =(\mathbf{a}+\mathbf{v})+\mathbf{b}=\mathbf{0}+\mathbf{b}=\mathbf{b}.
\]
Selanjutnya $0\mathbf{v}=(0+0)\mathbf{v}=0\mathbf{v}+0\mathbf{v}$. Menambahkan
$-(0\mathbf{v})$ pada kedua ruas memberi $\mathbf{0}=0\mathbf{v}$. Dari situ
$\mathbf{v}+(-1)\mathbf{v}=1\mathbf{v}+(-1)\mathbf{v}=(1-1)\mathbf{v}=0\mathbf{v}=\mathbf{0}$,
jadi $(-1)\mathbf{v}$ adalah invers $\mathbf{v}$, dan karena invers tunggal,
$(-1)\mathbf{v}=-\mathbf{v}$. Terakhir, jika $c\mathbf{v}=\mathbf{0}$ dan $c\neq 0$, maka
\[
  \mathbf{v}=1\mathbf{v}=(c^{-1}c)\mathbf{v}=c^{-1}(c\mathbf{v})=c^{-1}\mathbf{0}=\mathbf{0}.
\]
\hfill$\square$\par\medskip

Empat contoh berikut cukup untuk seluruh buku ini.

\emph{(i) Ruang koordinat.} $\mathbb{R}^{n}$ dengan penjumlahan dan perkalian skalar
koordinat demi koordinat. Di sini $\mathbf{0}=(0,\dots,0)$, dan kedelapan aksioma tinggal
sifat lapangan bilangan real yang sudah dibahas di Bab 3.

\emph{(ii) Ruang matriks.} $M_{2\times3}(\mathbb{R})$, himpunan matriks $2\times3$ berentri
real, dengan operasi pada setiap entri dan matriks nol sebagai $\mathbf{0}$. Tidak ada yang
berubah bila $2\times3$ diganti $m\times n$.

\emph{(iii) Ruang polinomial.} $P_{3}$ adalah himpunan polinomial berkoefisien real
berderajat paling tinggi $3$ beserta polinomial nol, dengan operasi suku demi suku. Jebakan
yang layak dicatat: polinomial berderajat \emph{tepat} $3$ tidak membentuk ruang vektor, sebab
$x^{3}$ dan $-x^{3}+x$ berderajat tiga tetapi jumlahnya $x$ berderajat satu. Kata ``paling
tinggi'' bukan hiasan.

\emph{(iv) Ruang fungsi.} Himpunan semua fungsi $\mathbb{R}\to\mathbb{R}$, ditulis
$\mathbb{R}^{\mathbb{R}}$, dengan $(f+g)(x)=f(x)+g(x)$ dan $(af)(x)=af(x)$. Fungsi nol adalah
$\mathbf{0}$, dan aksioma ruang vektor diwarisi dari $\mathbb{R}$ di setiap titik $x$.

Dari contoh (iv) muncul contoh yang lebih halus. Himpunan $C[0,1]$ fungsi kontinu pada
$[0,1]$ adalah ruang vektor, karena jumlah dan kelipatan skalar fungsi kontinu tetap kontinu.
Begitu pula barisan real yang konvergen, memakai sifat limit jumlah dari Bab 4. Keduanya
bagian dari ruang yang lebih besar, dan bagian semacam itu kita definisikan sekarang.

Kembali ke $\mathbb{R}^{2}$, tempat kedua operasi itu terlihat langsung.
Gambar~\ref{fig:ch06-dua-operasi} memakai $\mathbf{u}=(2,1)$ dan $\mathbf{v}=(1,2)$: jumlah
$\mathbf{u}+\mathbf{v}=(3,3)$ tercapai lewat dua urutan, dan $2\mathbf{u}$ serta
$-\mathbf{u}$ tetap segaris dengan $\mathbf{u}$.

\begin{figure}[htbp]
\centering
\begin{tikzpicture}[scale=0.85,>=Stealth,font=\small]
  \draw[->,layergrey] (-2.8,0) -- (5.0,0) node[below] {$x$};
  \draw[->,layergrey] (0,-1.8) -- (0,4.2) node[left] {$y$};
  \draw[dashed,layergrey] (2,1) -- (3,3);
  \draw[dashed,layergrey] (1,2) -- (3,3);
  \draw[->,thick,accent] (0,0) -- (2,1) node[below right] {$\mathbf{u}$};
  \draw[->,thick,layerindigo] (0,0) -- (1,2) node[above left] {$\mathbf{v}$};
  \draw[->,thick,layerteal] (0,0) -- (3,3) node[right] {$\mathbf{u}+\mathbf{v}$};
  \draw[->,layerorange] (0,0) -- (4,2) node[right] {$2\mathbf{u}$};
  \draw[->,layerbrown] (0,0) -- (-2,-1) node[left] {$-\mathbf{u}$};
  \fill[ink] (0,0) circle (1.2pt);
  \node[below left] at (0,0) {$\mathbf{0}$};
\end{tikzpicture}
\caption{Dua operasi di $\mathbb{R}^{2}$. Panah tebal menunjukkan $\mathbf{u}=(2,1)$ dan
$\mathbf{v}=(1,2)$; garis putus-putus melengkapi jajar genjang, sehingga
$\mathbf{u}+\mathbf{v}=(3,3)$ tercapai lewat dua urutan yang sama. Panah tipis menunjukkan
$2\mathbf{u}=(4,2)$ dan $-\mathbf{u}=(-2,-1)$, keduanya tetap pada satu garis dengan
$\mathbf{u}$.}
\label{fig:ch06-dua-operasi}
\end{figure}

\kamus{ruang vektor}{vector space}{kotak pencampur: isinya boleh dijumlahkan dan diperbesar, dan hasilnya tidak pernah keluar dari kotak.}
\begin{bacakeras}
Delapan aksioma, dibaca seperti daftar periksa. Dua suara, bergantian; suara kedua selalu
menjawab dengan tiga kata.

Suara pertama: Tukar tempat. Suara kedua: Hasilnya sama.
Suara pertama: Pindah kurung. Suara kedua: Hasilnya sama.
Suara pertama: Ada vektor nol. Suara kedua: Tidak menambah apa pun.
Suara pertama: Ada vektor lawan. Suara kedua: Kembali ke nol.
Suara pertama: Skalar menyebar ke dua suku. Suara kedua: Kiri dan kanan sama.
Suara pertama: Dua skalar menjumlah. Suara kedua: Satu vektor menanggung.
Suara pertama: Dua skalar mengali. Suara kedua: Satu kali saja.
Suara pertama: Skalar satu. Suara kedua: Diam, tidak berubah.

Dua baris terakhir berbunyi hampir sama padahal berbeda: yang satu menjumlah takaran, yang
lain mengalikan takaran. Di situlah kekeliruan hitung paling sering lahir.

Ucapkan sekali lagi, pelan, sebagai satu tarikan napas: tukar tempat, pindah kurung, ada nol,
ada lawan, menyebar, menjumlah, mengali, diam.
\end{bacakeras}


\section{Subruang dan uji subruang}
\tingkat{dasar}

Subruang adalah bagian dari ruang vektor yang tetap berperilaku sebagai ruang vektor walaupun
kita hanya memakai anggota bagian itu.

\begin{definisi}[Subruang]
\index{subruang}
Himpunan $W\subseteq V$ disebut \emph{subruang} dari $V$ jika $W$ merupakan ruang vektor
dengan operasi yang sama seperti pada $V$.
\end{definisi}

Definisi itu menuntut delapan aksioma, dan memeriksa delapan hal setiap kali terlalu mahal.
Enam di antaranya gratis: komutatif, asosiatif, distributif berlaku di $W$ karena berlaku di
$V$. Yang tersisa tiga hal, dan itulah isi teorema berikut.

Jika $W$ tidak kosong dan tertutup terhadap kedua operasi, maka $W$ otomatis subruang.

\begin{teorema}[Uji subruang]
\index{uji subruang}
Misalkan $V$ ruang vektor dan $W\subseteq V$ dengan $W\neq\varnothing$. Maka $W$ subruang
dari $V$ jika dan hanya jika $W$ tertutup terhadap penjumlahan dan perkalian skalar.
\end{teorema}

\par\medskip\noindent\textbf{Bukti.}\ 
Jika $W$ subruang, hasil kedua operasi pada anggota $W$ harus menjadi anggota $W$, jadi
kedua penutupan berlaku. Sebaliknya, andaikan $W$ tidak kosong dan tertutup terhadap kedua
operasi. Ambil $\mathbf{w}\in W$. Menurut penutupan skalar, $\mathbf{0}=0\mathbf{w}\in W$,
jadi $W$ memuat netral. Juga $-\mathbf{w}=(-1)\mathbf{w}\in W$, jadi setiap anggota $W$
punya invers di dalam $W$. Aksioma yang tersisa berlaku karena berlaku di $V$ dan operasi
pada $W$ hanyalah pembatasan operasi pada $V$.
\hfill$\square$\par\medskip

Dalam praktik kita memeriksa tiga hal: $\mathbf{0}\in W$ dan kedua ketertutupan. Ketiganya
perlu; tiga contoh berikut menunjukkan bahwa kegagalan satu syarat tidak dapat ditutupi
syarat lain.
\begin{resep}
\begin{enumerate}
  \item Tulis syarat keanggotaan $W$ sebagai persamaan, bukan sebagai gambar atau nama.
  \item Masukkan $\mathbf{v}=\mathbf{0}$ ke syarat itu. Bila gagal, berhenti: bukan subruang.
  \item Ambil $\mathbf{u},\mathbf{v}\in W$ sembarang, periksa $\mathbf{u}+\mathbf{v}$ memakai
        syarat yang sama, bukan memakai contoh angka.
  \item Ambil $c\in\mathbb{R}$ sembarang dan $\mathbf{u}\in W$, periksa $c\mathbf{u}$; sertakan
        $c<0$ dan $0<c<1$.
\end{enumerate}
Peringatan urutan: jangan menukar langkah 2 dengan langkah 3 dan 4. Himpunan yang tidak memuat
$\mathbf{0}$ tidak perlu diuji ketertutupannya, tetapi berhenti setelah langkah 2 juga keliru:
memuat $\mathbf{0}$ belum cukup, seperti ditunjukkan contoh-contoh berikut.
\end{resep}


\emph{Garis tak melalui titik asal.} $W=\{(x,y)\in\mathbb{R}^{2}:x+y=1\}$ tidak memuat
$\mathbf{0}$, jadi bukan subruang; ia juga tidak tertutup terhadap penjumlahan maupun
perkalian skalar, sebab $(1,0)+(0,1)$ dan $2(1,0)$ keduanya keluar dari $W$.

\emph{Gabungan dua sumbu.} $W=\{(x,0)\}\cup\{(0,y)\}$ memuat $\mathbf{0}$ dan tertutup
terhadap perkalian skalar, tetapi $(1,0)+(0,1)=(1,1)$ keluar dari $W$.

\emph{Kuadran pertama.} $W=\{(x,y):x\ge 0,\ y\ge 0\}$ tertutup terhadap penjumlahan,
tetapi $(-1)(1,1)=(-1,-1)$ keluar dari $W$.

Ketiga contoh gagal pada syarat berbeda; Gambar~\ref{fig:ch06-garis} memperlihatkan pasangan
paling sederhana, yaitu garis melalui titik asal dan garis yang melewatinya.

\begin{figure}[htbp]
\centering
\begin{tikzpicture}[scale=0.72,>=Stealth,font=\small]
  \draw[->,layergrey] (-1.6,0) -- (3.2,0) node[below] {$x$};
  \draw[->,layergrey] (0,-1.6) -- (0,3.2) node[left] {$y$};
  \draw[very thick,layerteal] (-1.4,-1.4) -- (2.8,2.8);
  \fill[ink] (0,0) circle (1.3pt) node[below left] {$\mathbf{0}$};
  \fill[accent] (2,2) circle (1.5pt) node[above right] {$(2,2)$};
  \node at (0.8,-2.4) {\emph{garis melalui titik asal}};
  \begin{scope}[xshift=7cm]
    \draw[->,layergrey] (-1.6,0) -- (3.2,0) node[below] {$x$};
    \draw[->,layergrey] (0,-1.6) -- (0,3.2) node[left] {$y$};
    \draw[very thick,accent] (-0.8,1.8) -- (2.4,-1.4);
    \fill[accent] (1,0) circle (1.5pt) node[below right] {$(1,0)$};
    \fill[accent] (0,1) circle (1.5pt) node[above left] {$(0,1)$};
    \draw[ink,thick] (0,0) circle (3pt);
    \node[below left] at (0,0) {$\mathbf{0}$};
    \node at (0.8,-2.4) {\emph{garis tanpa titik asal}};
  \end{scope}
\end{tikzpicture}
\caption{Dua himpunan di $\mathbb{R}^{2}$. Kiri: $\{(x,y):y=x\}$ memuat $\mathbf{0}$ dan
$(2,2)$, jadi tertutup terhadap kedua operasi. Kanan: $\{(x,y):x+y=1\}$ memuat $(1,0)$ dan
$(0,1)$, tetapi titik asal ditandai lingkaran kosong karena tidak termasuk; satu titik yang
hilang sudah cukup membatalkan subruang.}
\label{fig:ch06-garis}
\end{figure}



Sebaliknya, dua kelas contoh berikut selalu lolos.

\emph{Ruang nol.} Untuk matriks $A$ berukuran $m\times n$, himpunan
$\{\mathbf{x}:A\mathbf{x}=\mathbf{0}\}$ adalah subruang, sebab $A\mathbf{0}=\mathbf{0}$,
$A(\mathbf{x}+\mathbf{y})=A\mathbf{x}+A\mathbf{y}=\mathbf{0}$, dan
$A(c\mathbf{x})=cA\mathbf{x}=\mathbf{0}$. Himpunan itu disebut ruang nol atau kernel $A$,
ditulis $\ker A$, dan dibahas lanjut di Bab 7.

\emph{Polinomial yang lenyap di satu titik.} Himpunan
$\{p\in P_{3}:p(1)=0\}$ adalah subruang $P_{3}$, karena evaluasi di titik $1$ linear:
$(p+q)(1)=p(1)+q(1)$ dan $(cp)(1)=c\,p(1)$.

Kedua contoh itu sama-sama mempertahankan himpunan solusi persamaan homogen. Itu bukan
kebetulan.
\begin{tantangan}
Contoh kuadran pertama gagal karena satu skalar negatif. Carilah himpunan
$W\subseteq\mathbb{R}^{2}$ yang memuat $\mathbf{0}$, tertutup terhadap penjumlahan, tidak
tertutup terhadap perkalian skalar, dan bukan gabungan dua sumbu, bukan kuadran, bukan pula
himpunan berkoordinat bulat.

Petunjuk 1: jangan mengatur arah panahnya; atur bilangan apa yang boleh muncul pada jumlah
kedua koordinat.

Petunjuk 2: minta $x+y$ bulat. Maka $(1,0)+(0,1)=(1,1)$ tetap di dalam, sedangkan
$\tfrac{1}{2}(1,0)$ keluar.

Kalau mentok, kembali ke bagian ``Subruang dan uji subruang'' dan bandingkan dengan contoh
garis $x+y=1$; bedanya hanya pada ruas kanan persamaan.
\end{tantangan}


\kamus{subruang}{subspace}{lorong di dalam gedung: kita boleh berjalan bolak-balik di lorong itu, tetapi tidak pernah keluar darinya.}

\begin{untukanak}
Bayangkan dua lorong lurus di sebuah gedung yang bertemu di satu titik, seperti huruf X.
Setiap lorong, kalau kita berjalan di dalamnya, tetap merupakan lorong: kita bisa maju,
mundur, dan berhenti di mana saja tanpa meninggalkannya. Begitu juga kalau kita berjalan
dua kali lebih jauh di lorong yang sama; kita masih di lorong itu.

Sekarang perhatikan gabungan kedua lorong. Ambil satu langkah di lorong pertama, lalu satu
langkah di lorong kedua. Tempat kita berhenti tidak berada di lorong mana pun. Jadi gabungan
dua lorong bukan lorong, walaupun masing-masing lorong adalah lorong. Inilah hal yang paling
sering dilupakan tentang kata ``tertutup'': yang diperiksa bukan anggotanya satu per satu,
melainkan hasil mencampurnya.

Pertanyaan yang membuat penasaran: kalau ada tiga lorong bertemu di satu titik, apakah
gabungannya lorong? Kadang ya, kalau ketiganya lorong yang sama. Di bab ini kita akan tahu
mengapa ``tertutup terhadap pencampuran'' jauh lebih kuat daripada ``memuat titik temu'', dan
mengapa satu contoh langkah yang keluar sudah cukup membatalkan semuanya.
\end{untukanak}

\section{Kombinasi linear dan rentang}
\tingkat{menengah}

Kombinasi linear adalah hasil mencampur beberapa vektor dengan takaran tertentu. Adapun
rentang adalah himpunan semua hasil campuran itu, untuk segala takaran yang mungkin.

\begin{definisi}[Kombinasi linear dan rentang]
\index{kombinasi linear}
\index{rentang}
Vektor $\mathbf{v}$ adalah \emph{kombinasi linear} dari
$\mathbf{v}_{1},\dots,\mathbf{v}_{k}$ jika ada skalar $c_{1},\dots,c_{k}$ dengan
$\mathbf{v}=c_{1}\mathbf{v}_{1}+\dots+c_{k}\mathbf{v}_{k}$. \emph{Rentang} himpunan
$S\subseteq V$, ditulis $\operatorname{span}(S)$, adalah himpunan semua kombinasi linear
berhingga anggota $S$. Untuk $S=\varnothing$ kita tetapkan
$\operatorname{span}(S)=\{\mathbf{0}\}$.
\end{definisi}

Rentang selalu subruang, dan ia subruang terkecil yang memuat $S$. Pernyataan pertama
mengatakan rentang layak disebut ruang; yang kedua, tidak ada ruang lain yang lebih hemat.

\begin{lemma}
Untuk setiap $S\subseteq V$, $\operatorname{span}(S)$ adalah subruang $V$, dan setiap
subruang yang memuat $S$ juga memuat $\operatorname{span}(S)$.
\end{lemma}

\par\medskip\noindent\textbf{Bukti.}\ 
Untuk $S=\varnothing$ pernyataan benar karena $\{\mathbf{0}\}$ subruang. Untuk $S$ tak
kosong, ambil $\mathbf{u}=\sum a_{i}\mathbf{s}_{i}$ dan $\mathbf{v}=\sum b_{j}\mathbf{t}_{j}$
dua kombinasi linear berhingga anggota $S$. Maka
$\mathbf{u}+\mathbf{v}=\sum a_{i}\mathbf{s}_{i}+\sum b_{j}\mathbf{t}_{j}$ dan
$c\mathbf{u}=\sum (ca_{i})\mathbf{s}_{i}$ juga kombinasi linear berhingga anggota $S$, jadi
rentang tertutup terhadap kedua operasi; ia juga memuat $\mathbf{0}=0\mathbf{s}$ untuk
sembarang $\mathbf{s}\in S$. Jika $W$ subruang dan $S\subseteq W$, maka setiap
$\mathbf{s}\in S$ ada di $W$, dan ketertutupan $W$ memaksa setiap kombinasi linear berhingga
anggota $S$ juga ada di $W$.
\hfill$\square$\par\medskip

Karena itu $\operatorname{span}(S)$ adalah irisan semua subruang yang memuat $S$, dan
menambahkan vektor yang sudah di dalam rentang tidak mengubah apa pun:
$\operatorname{span}(\operatorname{span}(S))=\operatorname{span}(S)$.

Jika $\mathbf{b}$ berada di rentang kolom-kolom $A$, maka sistem $A\mathbf{x}=\mathbf{b}$
punya solusi.

\begin{proposisi}
Misalkan $A$ matriks $m\times n$ dengan kolom $\mathbf{a}_{1},\dots,\mathbf{a}_{n}$ dan
$\mathbf{b}\in\mathbb{R}^{m}$. Sistem $A\mathbf{x}=\mathbf{b}$ konsisten jika dan hanya jika
$\mathbf{b}\in\operatorname{span}(\mathbf{a}_{1},\dots,\mathbf{a}_{n})$.
\end{proposisi}

\par\medskip\noindent\textbf{Bukti.}\ 
Menulis $A\mathbf{x}=\mathbf{b}$ koordinat demi koordinat sama dengan menulis
$x_{1}\mathbf{a}_{1}+\dots+x_{n}\mathbf{a}_{n}=\mathbf{b}$. Jadi solusi $\mathbf{x}$ ada
tepat ketika $\mathbf{b}$ merupakan kombinasi linear kolom-kolom $A$ dengan koefisien
$x_{1},\dots,x_{n}$.
\hfill$\square$\par\medskip

Himpunan $\operatorname{span}(\mathbf{a}_{1},\dots,\mathbf{a}_{n})$ disebut
\emph{ruang kolom} $A$, dan ia subruang $\mathbb{R}^{m}$ menurut lemma di atas. Untuk tabel
pembuka, ruang kolom itu dibangun empat kolom, dan dua di antaranya berlebih.

Sebagai contoh kecil, di $\mathbb{R}^{3}$ tinjau $\mathbf{s}_{1}=(1,0,1)$ dan
$\mathbf{s}_{2}=(0,1,1)$. Setiap kombinasi linear berbentuk $(a,b,a+b)$, jadi rentangnya
bidang $z=x+y$. Vektor $(1,2,4)$ di luar bidang itu karena $4\neq 1+2$, sedangkan $(1,2,3)$
di dalamnya dengan $a=1$, $b=2$.

\begin{tangan}
Bahan: dua pensil sama panjang dan satu meja datar. Tandai satu titik di meja dengan ujung
jari; sebut itu titik asal.

Letakkan pensil pertama dengan ekor di titik asal. Letakkan pensil kedua dengan ekor di ujung
pensil pertama, menyiku dengan yang pertama. Sekarang geser pensil kedua tanpa memutarnya
sampai ekornya kembali ke titik asal. Ujungnya berpindah ke tempat lain, dan tempat itu adalah
jumlah kedua pensil: dua urutan penjumlahan berhenti di titik yang sama.

Sekarang skalarnya. Balik satu pensil, arahnya menjadi negatif. Sambung dua pensil yang sama
dari titik asal, panjangnya menjadi dua kali. Yang harus pembaca lihat: semua ujung itu tetap
pada satu garis lurus, sejauh apa pun pensil disambung. Satu pensil hanya menutupi satu garis
di meja; dua pensil yang tidak sejajar menutupi seluruh permukaan meja. Itu bedanya garis dan
bidang, dan tidak ada pensil ketiga yang menambah permukaan baru.
\end{tangan}



\section{Kebebasan linear dan ketergantungan}
\tingkat{menengah}

Himpunan vektor disebut bebas linear jika tidak ada satu pun anggotanya yang dapat dicampur
dari anggota lain. Ia bergantung linear jika ada anggota yang berlebih.

\begin{definisi}[Bebas dan bergantung linear]
\index{kebebasan linear}
Himpunan $\{\mathbf{v}_{1},\dots,\mathbf{v}_{k}\}\subseteq V$ \emph{bebas linear} jika
satu-satunya pilihan skalar yang memenuhi
$c_{1}\mathbf{v}_{1}+\dots+c_{k}\mathbf{v}_{k}=\mathbf{0}$ adalah
$c_{1}=\dots=c_{k}=0$. Jika ada pilihan lain, himpunan itu \emph{bergantung linear}.
\end{definisi}

Dua kesepakatan kecil: himpunan kosong bebas linear karena tidak ada koefisien yang perlu
diperiksa, dan $\{\mathbf{0}\}$ bergantung linear karena $1\cdot\mathbf{0}=\mathbf{0}$
adalah relasi nontrivial.

Dalam himpunan yang bergantung linear selalu ada satu vektor yang terletak di rentang
vektor-vektor sebelumnya.

\begin{lemma}[Ketergantungan]
Misalkan $\mathbf{v}_{1}\neq\mathbf{0}$ dan $S=\{\mathbf{v}_{1},\dots,\mathbf{v}_{k}\}$.
Himpunan $S$ bergantung linear jika dan hanya jika ada indeks $j\ge 2$ dengan
$\mathbf{v}_{j}\in\operatorname{span}(\mathbf{v}_{1},\dots,\mathbf{v}_{j-1})$.
\end{lemma}

\par\medskip\noindent\textbf{Bukti.}\ 
Jika $\mathbf{v}_{j}=\sum_{i<j}c_{i}\mathbf{v}_{i}$, maka
$\sum_{i<j}c_{i}\mathbf{v}_{i}-\mathbf{v}_{j}=\mathbf{0}$ adalah relasi nontrivial, karena
koefisien $\mathbf{v}_{j}$ adalah $-1$; jadi $S$ bergantung. Sebaliknya, jika $S$ bergantung
maka ada relasi nontrivial $\sum_{i=1}^{k}c_{i}\mathbf{v}_{i}=\mathbf{0}$. Pilih $j$ terbesar
dengan $c_{j}\neq 0$. Karena $\mathbf{v}_{1}\neq\mathbf{0}$, indeks $j$ bukan $1$, jadi
$j\ge 2$ dan
\[
  \mathbf{v}_{j}=-\frac{1}{c_{j}}\sum_{i<j}c_{i}\mathbf{v}_{i}
  \in\operatorname{span}(\mathbf{v}_{1},\dots,\mathbf{v}_{j-1}).
\]
\hfill$\square$\par\medskip

Tiga akibatnya kita pakai terus: himpunan yang memuat $\mathbf{0}$ selalu bergantung; dua
vektor bergantung tepat ketika salah satunya kelipatan yang lain; dan setiap himpunan bagian
dari himpunan bebas linear juga bebas linear.

Sebagai contoh, tinjau $\mathbf{v}_{1}=(1,2,-1)$, $\mathbf{v}_{2}=(0,1,3)$, dan
$\mathbf{v}_{3}=(2,5,1)$ di $\mathbb{R}^{3}$. Karena
$2\mathbf{v}_{1}+\mathbf{v}_{2}=(2,5,1)=\mathbf{v}_{3}$, relasi
$2\mathbf{v}_{1}+\mathbf{v}_{2}-\mathbf{v}_{3}=\mathbf{0}$ nontrivial dan ketiganya
bergantung. Tidak satu pun vektor itu nol, dan tidak ada dua yang berkelipatan; ketergantungan
baru muncul pada tiga vektor.

Gambar~\ref{fig:ch06-bidang} menaruh keduanya berdampingan: rentang dua vektor sebagai
bidang, dan tiga vektor yang salah satunya sudah di rentang dua lainnya.

\begin{figure}[htbp]
\centering
\begin{tikzpicture}[scale=0.8,>=Stealth,font=\small]
  \begin{scope}[x={(1cm,0cm)}, y={(0.40cm,0.28cm)}, z={(0cm,1cm)}]
    \fill[layerteal!10] (0,0,0) -- (2.4,0,2.4) -- (2.4,1.8,4.2) -- (0,1.8,2.4) -- cycle;
    \draw[layerteal] (0,0,0) -- (2.4,0,2.4) -- (2.4,1.8,4.2) -- (0,1.8,2.4) -- cycle;
    \draw[->,layergrey] (0,0,0) -- (3.0,0,0) node[below] {$x$};
    \draw[->,layergrey] (0,0,0) -- (0,2.2,0) node[right] {$y$};
    \draw[->,layergrey] (0,0,0) -- (0,0,2.6) node[left] {$z$};
    \draw[->,thick,accent] (0,0,0) -- (1,0,1)
      node[above left,font=\scriptsize,inner sep=1pt] {$\mathbf{s}_{1}$};
    \draw[->,thick,layerindigo] (0,0,0) -- (0,1,1)
      node[right,font=\scriptsize,inner sep=1pt] {$\mathbf{s}_{2}$};
    \fill[layerjade] (1,2,3) circle (1.6pt);
    \node[right,font=\scriptsize] at (1,2,3) {$(1,2,3)$};
    \draw[ink] (1,2,4) circle (2pt);
    \node[right,font=\scriptsize] at (1,2,4) {$(1,2,4)$};
  \end{scope}
  \begin{scope}[xshift=6.4cm]
    \draw[->,layergrey] (-0.4,0) -- (3.2,0) node[below] {$x$};
    \draw[->,layergrey] (0,-0.4) -- (0,2.6) node[left] {$y$};
    \draw[dashed,layergrey] (0,0) -- (2,1);
    \draw[dashed,layergrey] (2,1) -- (2.4,2);
    \draw[->,thick,accent] (0,0) -- (1,0.5) node[below right] {$\mathbf{v}_{1}$};
    \draw[->,thick,layerindigo] (0,0) -- (0.4,1) node[above left] {$\mathbf{v}_{2}$};
    \draw[->,very thick,layerorange] (0,0) -- (2.4,2) node[right] {$\mathbf{v}_{3}$};
    \fill[ink] (0,0) circle (1.2pt);
  \end{scope}
\end{tikzpicture}
\caption{Kiri: $\operatorname{span}\{(1,0,1),(0,1,1)\}$ adalah bidang $z=x+y$; titik
$(1,2,3)$ berada di dalamnya, titik $(1,2,4)$ di luarnya. Kanan: sketsa dua dimensi dari
$\mathbf{v}_{1}=(1,2,-1)$ dan $\mathbf{v}_{2}=(0,1,3)$ dengan
$\mathbf{v}_{3}=2\mathbf{v}_{1}+\mathbf{v}_{2}=(2,5,1)$; garis putus-putus menunjukkan
$2\mathbf{v}_{1}$ ditambah $\mathbf{v}_{2}$ mendarat tepat di ujung $\mathbf{v}_{3}$, sehingga
ketiganya bergantung.}
\label{fig:ch06-bidang}
\end{figure}



Di $\mathbb{R}^{2}$ setiap tiga vektor selalu bergantung, tanpa perlu menghitung apa pun.

\begin{proposisi}
Setiap himpunan tiga vektor di $\mathbb{R}^{2}$ bergantung linear.
\end{proposisi}

\par\medskip\noindent\textbf{Bukti.}\ 
Misalkan $S=\{\mathbf{u},\mathbf{v},\mathbf{w}\}\subseteq\mathbb{R}^{2}$. Jika dua di
antaranya berkelipatan, $S$ bergantung. Jika tidak, $\mathbf{u}$ dan $\mathbf{v}$ bebas
linear. Ambil $\mathbf{w}=(w_{1},w_{2})$ dan cari $a,b$ dengan
$a\mathbf{u}+b\mathbf{v}=\mathbf{w}$. Menuliskan komponennya memberi dua persamaan linear
dalam $a$ dan $b$, dan karena $\mathbf{u}$, $\mathbf{v}$ tidak berkelipatan, sistem itu punya
solusi tunggal. Jadi $\mathbf{w}\in\operatorname{span}(\mathbf{u},\mathbf{v})$.
\hfill$\square$\par\medskip

Di $\mathbb{R}^{3}$ pernyataan serupa berlaku untuk empat vektor, tetapi cara di atas akan
menjadi pengecekan kasus yang panjang; kita butuh argumen seragam.

Untuk memeriksa kebebasan secara hitung, susun vektor sebagai kolom matriks lalu selesaikan
$A\mathbf{c}=\mathbf{0}$; himpunan bebas linear tepat ketika sistem itu hanya punya solusi
trivial. Determinan memberi jalan pintas di $\mathbb{R}^{3}$, tetapi alat itu baru dibangun di
Bab 8.

\section{Basis dan dimensi}
\tingkat{lanjut}

Basis adalah daftar vektor yang cukup untuk membangun seluruh ruang dan tidak memuat satu pun
vektor berlebih. Kata ``cukup'' dan ``tidak berlebih'' persis dua syarat yang sudah kita punya.

\begin{definisi}[Basis dan dimensi]
\index{basis}
Himpunan $B\subseteq V$ disebut \emph{basis} dari $V$ jika $B$ bebas linear dan
$\operatorname{span}(B)=V$. Jika $V$ mempunyai basis berhingga, $V$ disebut
\emph{berdimensi hingga}; jika tidak, $V$ berdimensi tak hingga.
\end{definisi}

Definisi dimensi belum bisa ditulis, karena ia bergantung pada satu kenyataan yang belum
dibuktikan: dua basis dari ruang yang sama selalu berukuran sama.

Langkah pertama: penulisan dalam basis selalu tunggal.

\begin{lemma}[Ketunggalan koordinat]
Jika $B=\{\mathbf{b}_{1},\dots,\mathbf{b}_{n}\}$ basis $V$, maka setiap $\mathbf{v}\in V$
dapat ditulis secara tunggal sebagai $\mathbf{v}=a_{1}\mathbf{b}_{1}+\dots+a_{n}\mathbf{b}_{n}$.
\end{lemma}

\par\medskip\noindent\textbf{Bukti.}\ 
Keberadaan penulisan datang dari $\operatorname{span}(B)=V$. Andaikan ada dua penulisan
dengan koefisien $a_{i}$ dan $a_{i}'$. Pengurangan keduanya memberi
$\sum_{i}(a_{i}-a_{i}')\mathbf{b}_{i}=\mathbf{0}$. Karena $B$ bebas linear, semua koefisien
selisih itu nol, jadi $a_{i}=a_{i}'$ untuk setiap $i$.
\hfill$\square$\par\medskip

Angka $a_{1},\dots,a_{n}$ disebut koordinat $\mathbf{v}$ terhadap $B$; pertukaran basis
adalah pokok bahasan Bab 7, dan di sini kita hanya memerlukan ketunggalannya.

Langkah kedua: satu vektor dapat ditukar dengan vektor lain tanpa merusak sifat merentang.
Menukar satu vektor dalam himpunan merentang dengan kombinasi yang memuatnya tidak mengurangi
apa yang dapat dibangun.

\begin{lemma}[Penukaran]
\index{lemma penukaran}
Misalkan $\mathbf{u}=\sum_{\mathbf{s}\in S}a_{\mathbf{s}}\mathbf{s}$ jumlah berhingga dengan
$a_{\mathbf{s}_{0}}\neq 0$ untuk suatu $\mathbf{s}_{0}\in S$. Jika $S\cup\{\mathbf{u}\}$
merentang $V$, maka $(S\setminus\{\mathbf{s}_{0}\})\cup\{\mathbf{u}\}$ juga merentang $V$.
\end{lemma}

\par\medskip\noindent\textbf{Bukti.}\ 
Dari relasi itu,
$\mathbf{s}_{0}=a_{\mathbf{s}_{0}}^{-1}\bigl(\mathbf{u}-\sum_{\mathbf{s}\neq\mathbf{s}_{0}}a_{\mathbf{s}}\mathbf{s}\bigr)$,
jadi $\mathbf{s}_{0}$ adalah kombinasi linear anggota himpunan kecil
$T=(S\setminus\{\mathbf{s}_{0}\})\cup\{\mathbf{u}\}$. Setiap anggota $S\cup\{\mathbf{u}\}$
karena itu berada di $\operatorname{span}(T)$, sehingga
$\operatorname{span}(S\cup\{\mathbf{u}\})\subseteq\operatorname{span}(T)$. Arah sebaliknya
jelas karena $T\subseteq S\cup\{\mathbf{u}\}$. Jadi kedua rentang sama dengan $V$.
\hfill$\square$\par\medskip

Langkah ketiga adalah inti bab ini.

Jika ruang dapat dibangun oleh $n$ vektor, maka mustahil menemukan $n+1$ vektor yang bebas
linear di dalamnya.

\begin{teorema}[Pertukaran]
\index{teorema pertukaran}
Jika $V$ direntang oleh $n$ vektor dan $\mathbf{w}_{1},\dots,\mathbf{w}_{m}$ bebas linear di
$V$, maka $m\le n$.
\end{teorema}

\par\medskip\noindent\textbf{Bukti.}\ 
Misalkan $S_{0}=\{\mathbf{v}_{1},\dots,\mathbf{v}_{n}\}$ merentang $V$. Untuk $j=1$:
$\mathbf{w}_{1}\neq\mathbf{0}$ adalah kombinasi linear anggota $S_{0}$ dengan sekurangnya
satu koefisien tak nol, sehingga lemma penukaran memberi $T_{1}\subseteq S_{0}$ berukuran
$n-1$ dengan $\{\mathbf{w}_{1}\}\cup T_{1}$ merentang.

Andaikan $\{\mathbf{w}_{1},\dots,\mathbf{w}_{j}\}\cup T_{j}$ merentang dengan
$|T_{j}|=n-j$. Vektor $\mathbf{w}_{j+1}$ adalah kombinasi linear himpunan itu. Jika semua
koefisien pada $T_{j}$ nol, maka $\mathbf{w}_{j+1}$ menjadi kombinasi linear
$\mathbf{w}_{1},\dots,\mathbf{w}_{j}$, yang mustahil karena himpunan itu bebas linear. Jadi
ada anggota $T_{j}$ berkoefisien tak nol, dan lemma penukaran memberi $T_{j+1}$ berukuran
$n-j-1$ yang tetap merentang bersama $\mathbf{w}_{1},\dots,\mathbf{w}_{j+1}$.

Setelah $m$ langkah kita memperoleh himpunan merentang berukuran $n-m$; ukuran himpunan tidak
boleh negatif, jadi $m\le n$.
\hfill$\square$\par\medskip

Sekarang ukuran sebuah ruang dapat ditulis: banyaknya vektor dalam daftar paling hematnya,
dan angka itu tidak bergantung pada daftar yang dipilih.

Dua basis dari ruang berdimensi hingga selalu sama banyak.

\begin{akibat}
Jika $V$ berdimensi hingga dengan basis $B_{1}$ berukuran $n$ dan basis $B_{2}$ berukuran
$m$, maka $n=m$.
\end{akibat}

\par\medskip\noindent\textbf{Bukti.}\ 
$B_{1}$ merentang dan $B_{2}$ bebas linear, jadi teorema pertukaran memberi $m\le n$.
Sebaliknya $B_{2}$ merentang dan $B_{1}$ bebas linear, sehingga $n\le m$.
\hfill$\square$\par\medskip

Banyaknya vektor dalam daftar paling hemat itu cukup penting untuk diberi nama sendiri.

\begin{definisi}[Dimensi]
\index{dimensi}
Untuk $V$ berdimensi hingga, $\dim V$ adalah banyaknya anggota sembarang basis $V$. Ruang
nol $\{\mathbf{0}\}$ ditetapkan berdimensi $0$.
\end{definisi}

Ukuran subruang tidak pernah melebihi ukuran ruangnya; daftar yang panjangnya pas dengan
dimensi langsung menjadi basis; dan daftar yang bebas linear selalu bisa dilengkapi sampai
penuh.

\begin{akibat}
Misalkan $\dim V=n$. (a) Setiap himpunan $n$ vektor yang bebas linear adalah basis.
(b) Setiap himpunan $n$ vektor yang merentang adalah basis. (c) Setiap himpunan bebas linear
dapat diperluas menjadi basis. (d) Jika $U$ subruang $V$, maka $\dim U\le\dim V$, dan
kesamaan berlaku hanya bila $U=V$.
\end{akibat}

\par\medskip\noindent\textbf{Bukti.}\ 
Untuk (a), himpunan bebas linear berukuran $n$ yang tidak merentang dapat ditambah satu
vektor di luar rentangnya, dan hasilnya tetap bebas linear; jadi ada $n+1$ vektor bebas
linear, dan teorema pertukaran melawannya. Untuk (b), himpunan merentang berukuran $n$ yang
bergantung dapat kehilangan satu vektor lewat lemma ketergantungan dan tetap merentang dengan
$n-1$ vektor; lagi-lagi bertentangan dengan pertukaran. Untuk (c), tambahkan vektor di luar
rentang selama masih ada; proses berhenti setelah paling banyak $n$ langkah, dan saat berhenti
rentangnya sudah seluruh $V$. Untuk (d), basis $U$ adalah himpunan bebas linear di $V$, jadi
ukurannya paling banyak $n$; bila ukurannya $n$, bagian (a) membuatnya basis $V$, sehingga
$U=V$.
\hfill$\square$\par\medskip

Dengan demikian $\dim\mathbb{R}^{n}=n$, $\dim P_{n}=n+1$ dengan basis $1,x,\dots,x^{n}$,
$\dim M_{m\times n}(\mathbb{R})=mn$ dengan basis matriks yang tepat satu entrinya $1$, dan
$\dim\{\mathbf{0}\}=0$. Untuk ruang fungsi seperti $\mathbb{R}^{\mathbb{R}}$ tidak ada basis
berhingga; buktinya memerlukan aksioma pilihan dan di luar cakupan buku ini.

\kamus{basis}{basis}{daftar kaleng cat paling hemat: cukup untuk membuat semua warna, dan tidak ada satu kaleng pun yang bisa dibuang.}

\kamus{dimensi}{dimension}{banyaknya kaleng dalam daftar paling hemat itu; angkanya milik ruang, bukan milik daftar yang kebetulan kita pilih.}

\section{Jumlah langsung dan subruang komplemen}
\tingkat{lanjut}

Jumlah dua subruang adalah himpunan semua hasil menambah satu anggota dari masing-masing
subruang.

\begin{definisi}[Jumlah subruang]
\index{jumlah subruang}
Untuk subruang $U,W\subseteq V$, \emph{jumlah} $U+W$ adalah
$\{\mathbf{u}+\mathbf{w}:\mathbf{u}\in U,\ \mathbf{w}\in W\}$.
\end{definisi}

Ia berbeda dari gabungan, yang hampir selalu bukan subruang seperti contoh dua sumbu.

Jumlahnya sendiri sebuah subruang, dan tidak ada subruang lain yang lebih kecil yang
memuat keduanya.

\begin{proposisi}
Untuk subruang $U,W\subseteq V$, himpunan $U+W$ adalah subruang $V$, dan ia termuat di
setiap subruang yang memuat $U\cup W$.
\end{proposisi}

\par\medskip\noindent\textbf{Bukti.}\ 
Himpunan $U+W$ memuat $\mathbf{0}=\mathbf{0}+\mathbf{0}$. Jika
$\mathbf{u}_{1}+\mathbf{w}_{1}$ dan $\mathbf{u}_{2}+\mathbf{w}_{2}$ dua anggotanya, jumlahnya
$(\mathbf{u}_{1}+\mathbf{u}_{2})+(\mathbf{w}_{1}+\mathbf{w}_{2})$ dengan suku pertama di $U$
dan kedua di $W$, dan $c(\mathbf{u}_{1}+\mathbf{w}_{1})=c\mathbf{u}_{1}+c\mathbf{w}_{1}$ juga
demikian; jadi $U+W$ subruang. Jika $Z$ subruang memuat $U$ dan $W$, maka $Z$ memuat setiap
$\mathbf{u}+\mathbf{w}$ karena tertutup terhadap penjumlahan, jadi $U+W\subseteq Z$.
\hfill$\square$\par\medskip

Ukuran jumlah bukan jumlah ukuran, karena bagian yang beririsan terhitung dua kali.

$\dim(U+W)$ sama dengan $\dim U$ ditambah $\dim W$ dikurangi dimensi irisan.

\begin{teorema}[Dimensi jumlah]
\index{dimensi jumlah}
Jika $U$ dan $W$ subruang ruang berdimensi hingga $V$, maka
\[
  \dim(U+W)=\dim U+\dim W-\dim(U\cap W).
\]
\end{teorema}

\par\medskip\noindent\textbf{Bukti.}\ 
Irisan $U\cap W$ subruang, jadi berdimensi hingga. Ambil basis
$\{\mathbf{z}_{1},\dots,\mathbf{z}_{m}\}$ dari $U\cap W$. Karena bebas linear di $U$,
himpunan itu diperluas menjadi basis $U$ dengan menambahkan $\mathbf{u}_{1},\dots,
\mathbf{u}_{j}$, sehingga $\dim U=m+j$; dengan cara sama, menambahkan
$\mathbf{w}_{1},\dots,\mathbf{w}_{k}$ menghasilkan basis $W$, sehingga $\dim W=m+k$.

Klaim: $\{\mathbf{z}_{1},\dots,\mathbf{z}_{m},\mathbf{u}_{1},\dots,\mathbf{u}_{j},
\mathbf{w}_{1},\dots,\mathbf{w}_{k}\}$ basis $U+W$. Himpunan itu memuat basis $U$ dan basis
$W$, jadi rentangnya memuat $U+W$; arah sebaliknya jelas. Untuk kebebasan, andaikan
\[
  \sum_{i}a_{i}\mathbf{z}_{i}+\sum_{r}b_{r}\mathbf{u}_{r}
  +\sum_{s}c_{s}\mathbf{w}_{s}=\mathbf{0}.
\]
Maka $\sum_{s}c_{s}\mathbf{w}_{s}=-\sum_{i}a_{i}\mathbf{z}_{i}-\sum_{r}b_{r}\mathbf{u}_{r}$
berada di $U$ sekaligus di $W$; jadi ia anggota $U\cap W$ dan dapat ditulis
$\sum_{i}d_{i}\mathbf{z}_{i}$. Maka $\sum_{s}c_{s}\mathbf{w}_{s}-\sum_{i}d_{i}\mathbf{z}_{i}
=\mathbf{0}$ adalah relasi linear dalam basis $W$, sehingga semua $c_{s}$ dan $d_{i}$ nol;
persamaan awal menyusut menjadi relasi linear dalam basis $U$, sehingga $a_{i}$ dan $b_{r}$
juga nol.

Jadi $\dim(U+W)=m+j+k=(m+j)+(m+k)-m=\dim U+\dim W-\dim(U\cap W)$.
\hfill$\square$\par\medskip

Sekarang pertanyaan kedua pembuka bab: kapan penguraian menjadi dua suku itu tunggal?

Jumlah $U+W$ disebut \emph{jumlah langsung}, ditulis $U\oplus W$, jika setiap anggota
$U+W$ dapat ditulis sebagai $\mathbf{u}+\mathbf{w}$ dengan cara tepat satu.

\begin{definisi}[Jumlah langsung]
\index{jumlah langsung}
Jika setiap $\mathbf{x}\in U+W$ memiliki tepat satu penulisan
$\mathbf{x}=\mathbf{u}+\mathbf{w}$ dengan $\mathbf{u}\in U$ dan $\mathbf{w}\in W$, maka
$U+W$ disebut jumlah langsung dan ditulis $U\oplus W$.
\end{definisi}

Ketunggalan seluruh penulisan dapat diperiksa hanya dengan melihat penulisan nol.

$U+W$ jumlah langsung tepat ketika $U\cap W=\{\mathbf{0}\}$.

\begin{teorema}[Kriteria jumlah langsung]
Misalkan $U$ dan $W$ subruang $V$. Maka $U+W=U\oplus W$ jika dan hanya jika
$U\cap W=\{\mathbf{0}\}$.
\end{teorema}

\par\medskip\noindent\textbf{Bukti.}\ 
Andaikan $U+W$ jumlah langsung dan ambil $\mathbf{x}\in U\cap W$. Maka
$\mathbf{0}=\mathbf{x}+(-\mathbf{x})$ dengan $\mathbf{x}\in U$ dan $-\mathbf{x}\in W$, dan
juga $\mathbf{0}=\mathbf{0}+\mathbf{0}$. Ketunggalan memberi $\mathbf{x}=\mathbf{0}$.

Sebaliknya andaikan $U\cap W=\{\mathbf{0}\}$. Ambil $\mathbf{x}\in U+W$ dengan dua penulisan
$\mathbf{x}=\mathbf{u}_{1}+\mathbf{w}_{1}=\mathbf{u}_{2}+\mathbf{w}_{2}$. Maka
$\mathbf{u}_{1}-\mathbf{u}_{2}=\mathbf{w}_{2}-\mathbf{w}_{1}$ berada di $U$ dan di $W$,
jadi sama dengan $\mathbf{0}$. Akibatnya $\mathbf{u}_{1}=\mathbf{u}_{2}$ dan
$\mathbf{w}_{1}=\mathbf{w}_{2}$.
\hfill$\square$\par\medskip

Bagi dua subruang kriteria itu cukup. Bagi tiga subruang atau lebih, memeriksa irisan
berpasangan tidak cukup, dan contoh berikut menunjukkan mengapa.

Di $\mathbb{R}^{3}$ ambil $V_{1}=\operatorname{span}\{(1,0,0)\}$,
$V_{2}=\operatorname{span}\{(0,1,0)\}$, dan $V_{3}=\operatorname{span}\{(1,1,0)\}$. Setiap
irisan berpasangan hanya memuat $\mathbf{0}$, tetapi $V_{1}+V_{2}+V_{3}$ bukan jumlah
langsung: $(1,1,0)$ dapat ditulis $(1,0,0)+(0,1,0)+\mathbf{0}$ dan juga
$\mathbf{0}+\mathbf{0}+(1,1,0)$. Nol pun punya dua penulisan, yaitu
$\mathbf{0}+\mathbf{0}+\mathbf{0}$ dan $(1,0,0)+(0,1,0)+(-1,-1,0)$.

Untuk banyak suku, kriterianya kembali ke definisi: $\mathbf{0}=\mathbf{v}_{1}+\dots+
\mathbf{v}_{k}$ dengan $\mathbf{v}_{i}\in V_{i}$ hanya mungkin bila semua $\mathbf{v}_{i}$
nol.

Terakhir, setiap subruang dapat dipasangkan dengan subruang lain sehingga keduanya
menguraikan seluruh ruang.

\begin{teorema}[Komplemen]
\index{komplemen}
Jika $V$ berdimensi hingga dan $U$ subruang $V$, maka ada subruang $W\subseteq V$ dengan
$V=U\oplus W$. Dalam hal itu $\dim V=\dim U+\dim W$.
\end{teorema}

\par\medskip\noindent\textbf{Bukti.}\ 
Jika $U=\{\mathbf{0}\}$, ambil $W=V$. Jika tidak, ambil basis
$\{\mathbf{u}_{1},\dots,\mathbf{u}_{m}\}$ dari $U$ dan perluas menjadi basis
$\{\mathbf{u}_{1},\dots,\mathbf{u}_{m},\mathbf{w}_{1},\dots,\mathbf{w}_{k}\}$ dari $V$;
perluasan ini dimungkinkan oleh akibat (c) di bagian sebelumnya. Tetapkan
$W=\operatorname{span}\{\mathbf{w}_{1},\dots,\mathbf{w}_{k}\}$.

Setiap $\mathbf{x}\in V$ adalah kombinasi linear seluruh basis, dan memisahkan suku
$\mathbf{u}$ dari suku $\mathbf{w}$ menunjukkannya sebagai anggota $U+W$. Untuk irisan, ambil
$\mathbf{x}\in U\cap W$. Maka
$\mathbf{x}=\sum a_{i}\mathbf{u}_{i}=\sum b_{s}\mathbf{w}_{s}$, sehingga
$\sum a_{i}\mathbf{u}_{i}-\sum b_{s}\mathbf{w}_{s}=\mathbf{0}$. Ini relasi linear dalam basis
$V$, jadi semua koefisiennya nol, khususnya $\mathbf{x}=\mathbf{0}$. Menurut kriteria di
atas, $V=U\oplus W$, dan dari teorema dimensi jumlah,
$\dim V=\dim U+\dim W$.

Komplemen itu tidak tunggal. Di $\mathbb{R}^{2}$ dengan $U$ sumbu-$x$, setiap garis melalui
titik asal selain sumbu-$x$ adalah komplemen $U$; yang tunggal bukan pilihannya, melainkan
dimensinya. Di Bab 10, komplemen ortogonal membuat penguraian itu sekaligus terpendek.

Sekarang pertanyaan pembuka dapat dijawab. Untuk $\mathbf{w}=(1,1,1)$, misalkan
$U=\{\boldsymbol{\delta}:\mathbf{w}\cdot\boldsymbol{\delta}=0\}$ dan
$W=\operatorname{span}\{\mathbf{w}\}$. Irisannya hanya $\{\mathbf{0}\}$, sebab
$\mathbf{w}\cdot(t\mathbf{w})=3t$ yang nol hanya bila $t=0$; dengan $\dim U=2$ dan
$\dim W=1$ kita peroleh $U\oplus W=\mathbb{R}^{3}$. Setiap perturbasi terpecah tunggal menjadi
bagian yang tidak mengubah skor, yaitu anggota $U$, dan bagian yang mengubahnya, anggota $W$.

\kamus{jumlah langsung}{direct sum}{dua lorong yang hanya bertemu di titik asal, sehingga setiap tempat di gedung dapat dicapai dengan satu jalan yang tidak bisa diperselisihkan.}

\section{Contoh hitung}
\tingkat{menengah}

Empat hitungan berikut memakai seluruh isi bab; semuanya diperiksa ulang dengan Python di
bagian lab.

\subsection*{(a) Rank dan basis ruang kolom}

Kembali ke tabel pembuka. Eliminasi baris: $R_{2}\leftarrow R_{2}-2R_{1}$ mengubah baris
kedua menjadi $(0,0,1,1)$, dan $R_{3}\leftarrow R_{3}-3R_{1}$ mengubah baris ketiga menjadi
$(0,0,2,2)$. Lalu $R_{3}\leftarrow R_{3}-2R_{2}$ menghabisi baris ketiga, sehingga tersisa
$\mathbf{r}_{1}=(1,2,0,1)$ dan $(0,0,1,1)$.

Bentuk itu sudah eselon tereduksi, dengan kolom pivot $1$ dan $3$. Jadi
$\operatorname{rank}A=2$ dan ruang kolomnya berdimensi $2$ dengan basis
$\{\mathbf{c}_{1},\mathbf{c}_{3}\}=\{(1,2,3),(0,1,2)\}$.

Pemeriksaan langsung terhadap kolom yang dibuang:
\[
  \mathbf{c}_{2}=(2,4,6)=2\mathbf{c}_{1},
  \qquad
  \mathbf{c}_{4}=(1,3,5)=\mathbf{c}_{1}+\mathbf{c}_{3}.
\]
Keduanya eksak, tanpa pembulatan. Inilah jawaban ``dua'' untuk pertanyaan pembuka: kolom pivot
memberi himpunan bebas linear yang merentang ruang kolom, dan kolom tanpa pivot adalah
kombinasi kolom pivot di kirinya.

Ruang nol $A$ juga terbaca dari bentuk tereduksi. Dari $x_{1}+2x_{2}+x_{4}=0$ dan
$x_{3}+x_{4}=0$ kita peroleh $x_{1}=-2x_{2}-x_{4}$ dan $x_{3}=-x_{4}$, sehingga
\[
  \mathbf{x}=x_{2}\,(-2,1,0,0)+x_{4}\,(-1,0,-1,1),
\]
ditulis sebagai vektor baris agar ringkas. Jadi $\dim\ker A=2$, dan
$\operatorname{rank}A+\dim\ker A=2+2=4$ sama dengan banyak kolom; pernyataan umumnya adalah
rank-nullity, dibuktikan di Bab 7.

\subsection*{(b) Tiga vektor yang bergantung}

Tinjau $\mathbf{v}_{1}=(1,2,-1)$, $\mathbf{v}_{2}=(0,1,3)$, $\mathbf{v}_{3}=(2,5,1)$ sebagai
kolom $V$. Ekspansi kofaktor sepanjang baris pertama:
\[
  \det V=1(1\cdot1-3\cdot5)+0+2\bigl(2\cdot3-1\cdot(-1)\bigr)=(1-15)+2(6+1)=0.
\]
Determinan nol berarti kolomnya bergantung; relasinya ditebak lalu diperiksa:
$2\mathbf{v}_{1}+\mathbf{v}_{2}=(2,5,1)=\mathbf{v}_{3}$, jadi
$2\mathbf{v}_{1}+\mathbf{v}_{2}-\mathbf{v}_{3}=\mathbf{0}$ dan $\operatorname{rank}V=2$.

\subsection*{(c) Dimensi ruang polinomial}

$\dim P_{3}=4$, karena $\{1,x,x^{2},x^{3}\}$ basis. Himpunan
$B=\{1,\ 1+x,\ 1+x+x^{2},\ 1+x+x^{2}+x^{3}\}$ juga basis: matriks koordinatnya terhadap
basis standar berbentuk segitiga atas dengan diagonal $1$, jadi $B$ bebas linear, dan karena
jumlahnya $4=\dim P_{3}$, akibat (a) membuatnya basis.

Koordinat $p=2-x+3x^{3}$ terhadap $B$ dihitung dari derajat tertinggi ke bawah: $a_{4}=3$,
lalu $a_{3}+a_{4}=0$ memberi $a_{3}=-3$, $a_{2}+a_{3}+a_{4}=-1$ memberi $a_{2}=-1$, dan
konstanta memberi $a_{1}=3$. Jadi $p=3-(1+x)-3(1+x+x^{2})+3(1+x+x^{2}+x^{3})$, yang bila
diuraikan memberi $2-x+3x^{3}$.

\subsection*{(d) Penguraian jumlah langsung}

Ambil $U=\{(x,y,z):x+y+z=0\}$ dan $W=\operatorname{span}\{(1,1,1)\}$ di $\mathbb{R}^{3}$.
Basis $U$ adalah $\{(1,-1,0),(1,0,-1)\}$, jadi $\dim U=2$; dengan $\dim W=1$ dan
$U\cap W=\{\mathbf{0}\}$ kita peroleh $\mathbb{R}^{3}=U\oplus W$.

Uraikan $\mathbf{v}=(3,0,1)$ sebagai $\mathbf{u}+t(1,1,1)$. Syarat jumlah koordinat
$\mathbf{u}$ nol memberi $3+0+1-3t=0$, jadi $t=4/3$. Maka
\[
  \mathbf{v}=\underbrace{\Bigl(\tfrac53,-\tfrac43,-\tfrac13\Bigr)}_{\in U}
  +\underbrace{\Bigl(\tfrac43,\tfrac43,\tfrac43\Bigr)}_{\in W},
\]
dan $\tfrac53-\tfrac43-\tfrac13=0$ memastikan suku pertama memang di $U$. Menurut teorema
kriteria, tidak ada penguraian lain.

\kertas{Ulangi (a) tanpa melihat: kerjakan eliminasi pada $A$, tandai kolom pivot, lalu
tuliskan $\mathbf{c}_{2}$ dan $\mathbf{c}_{4}$ sebagai kombinasi $\mathbf{c}_{1}$ dan
$\mathbf{c}_{3}$. Kemudian hitung $\dim\ker A$ dari bentuk tereduksi dan periksa bahwa
jumlahnya dengan rank sama dengan banyak kolom. Terakhir, untuk matriks $V$ pada (b),
gantilah entri $(3,1)$ menjadi $0$ dan tentukan apakah kolomnya masih bergantung.}

\section{Lab komputasi}
\tingkat{lanjut}

Kode berikut dijalankan apa adanya dengan numpy. Fungsi \texttt{rref} ditulis sendiri supaya
kolom pivot terlihat; baris terakhirnya membersihkan sisa pembulatan dan tanda nol negatif.

\begin{lab}
\begin{verbatim}
import numpy as np

np.set_printoptions(precision=3, suppress=True)

def rref(M):
    M = M.astype(float).copy()
    r = 0
    for c in range(M.shape[1]):
        i = np.argmax(np.abs(M[r:, c])) + r
        if abs(M[i, c]) < 1e-12:
            continue
        M[[r, i]] = M[[i, r]]
        M[r] = M[r] / M[r, c]
        for j in range(M.shape[0]):
            if j != r:
                M[j] = M[j] - M[j, c] * M[r]
        r = r + 1
        if r == M.shape[0]:
            break
    M[np.abs(M) < 1e-12] = 0.0
    return M

A = np.array([[1, 2, 0, 1],
              [2, 4, 1, 3],
              [3, 6, 2, 5]], float)
print("RREF(A):"); print(rref(A))
print("rank(A) =", np.linalg.matrix_rank(A))
print("k2-2k1 =", A[:, 1] - 2 * A[:, 0])
print("k4-k1-k3 =", A[:, 3] - A[:, 0] - A[:, 2])

V = np.array([[1, 0, 2],
              [2, 1, 5],
              [-1, 3, 1]], float)
print("det(V) =", np.linalg.det(V))
print("2v1+v2-v3 =", 2 * V[:, 0] + V[:, 1] - V[:, 2])

x = np.array([0., 1, 2, 3])
M = np.vander(x, 4, increasing=True)
print("det(M) =", np.linalg.det(M))
print("rank(M) =", np.linalg.matrix_rank(M))

B = np.array([[1, 1, 1, 1],
              [0, 1, 1, 1],
              [0, 0, 1, 1],
              [0, 0, 0, 1]], float)
p = np.array([2, -1, 0, 3], float)
k = np.linalg.solve(B, p)
print("koordinat p =", k)
print("sisa =", B @ k - p)
\end{verbatim}
\end{lab}

Keluaran yang tercetak:
\begin{verbatim}
RREF(A):
[[1. 2. 0. 1.]
 [0. 0. 1. 1.]
 [0. 0. 0. 0.]]
rank(A) = 2
k2-2k1 = [0. 0. 0.]
k4-k1-k3 = [0. 0. 0.]
det(V) = 0.0
2v1+v2-v3 = [0. 0. 0.]
det(M) = 12.0
rank(M) = 4
koordinat p = [ 3. -1. -3.  3.]
sisa = [0. 0. 0. 0.]
\end{verbatim}

Empat hal terbaca dari keluaran itu. Kolom pivot $1$ dan $3$ serta
$\operatorname{rank}A=2$ cocok dengan hitungan tangan. Selisih
$\mathbf{c}_{2}-2\mathbf{c}_{1}$ dan $\mathbf{c}_{4}-\mathbf{c}_{1}-\mathbf{c}_{3}$ nol
eksak; pada data nyata selisih itu kecil tetapi tak nol, dan di situlah batas penerapan rank
secara numerik. $\det V=0$ menegaskan ketergantungan di (b), dan matriks Vandermonde titik $0,1,2,3$
berdeterminan $12$ serta berank $4$. Koordinat $p$ tercetak $(3,-1,-3,3)$, cocok dengan
hitungan tangan.

\begin{untukdosen}
Mulai dari tabel pembuka, bukan dari aksioma: minta mahasiswa menebak berapa kolom yang
mandiri, lalu tunjukkan bahwa tebakan ``empat'' salah. Aksioma masuk setelah itu, sebagai
jawaban atas kebutuhan. Uji subruang dan rentang pada pertemuan kedua, kebebasan linear dan
basis pada pertemuan ketiga, jumlah langsung pada pertemuan keempat.

Dua miskonsepsi khas. Pertama, memeriksa subruang hanya dengan ``apakah memuat nol'' lalu
berhenti; kuadran pertama memuat nol dan tertutup terhadap penjumlahan, tetapi gagal pada
kelipatan negatif. Kedua, menganggap basis berarti basis standar, sehingga $(5,7)$ dan
$(4,3)$ dianggap bukan basis $\mathbb{R}^{2}$; tekankan bahwa basis adalah sifat himpunan.

Aktivitas kelas: setiap kelompok menulis tiga vektor di $\mathbb{R}^{3}$ dan memutuskan bebas
atau bergantung lewat eliminasi; kelompok lain mengubah satu entri dan meminta keputusan
ulang. Mahasiswa siap lanjut bila dapat menjelaskan mengapa tiga vektor di $\mathbb{R}^{2}$
selalu bergantung.
\end{untukdosen}

\section{Ringkasan}
\tingkat{dasar}

\begin{ringkasan}
\begin{itemize}
  \item Ruang vektor: himpunan dengan penjumlahan dan perkalian skalar yang memenuhi delapan
        aksioma; anggotanya boleh vektor, matriks, polinomial, atau fungsi.
  \item Subruang cukup diperiksa dengan tiga hal: memuat $\mathbf{0}$ dan tertutup terhadap
        kedua operasi.
  \item $\operatorname{span}(S)$ adalah subruang terkecil yang memuat $S$; $\mathbf{b}$ di
        rentang kolom $A$ tepat ketika $A\mathbf{x}=\mathbf{b}$ konsisten.
  \item Bergantung linear: satu anggota berada di rentang anggota sebelumnya; tiga vektor di
        $\mathbb{R}^{2}$ selalu bergantung.
  \item Teorema pertukaran menutup $n+1$ vektor bebas linear di ruang berrentang $n$ vektor,
        sehingga semua basis sama ukuran dan $\dim V$ terdefinisi.
  \item $\dim(U+W)=\dim U+\dim W-\dim(U\cap W)$; $U+W$ langsung tepat ketika
        $U\cap W=\{\mathbf{0}\}$; setiap subruang berdimensi hingga punya komplemen.
\end{itemize}
\end{ringkasan}

\section{Latihan}
\tingkat{dasar}

\begin{latihan}
\begin{enumerate}
  \item Tentukan apakah $W=\{(x,y,z)\in\mathbb{R}^{3}:x-2y+z=0\ \text{dan}\ x+y=0\}$
        subruang $\mathbb{R}^{3}$. Bila ya, tentukan basis dan dimensinya.
  \item Buktikan bahwa himpunan semua matriks $2\times2$ yang entri diagonal utamanya nol
        adalah subruang $M_{2\times2}(\mathbb{R})$, lalu tentukan basis dan dimensinya.
  \item Tunjukkan bahwa $S=\{1+x,\ 1+x^{2},\ 2+x+x^{2}\}$ bergantung linear di $P_{2}$.
        Tentukan $\dim\operatorname{span}(S)$, lalu perluas himpunan bagian bebas linear
        terbesar dari $S$ menjadi basis $P_{2}$.
  \item Buktikan: jika $\{\mathbf{v}_{1},\mathbf{v}_{2},\mathbf{v}_{3}\}$ bebas linear di
        $V$, maka $\{\mathbf{v}_{1}+\mathbf{v}_{2},\ \mathbf{v}_{2}+\mathbf{v}_{3},\
        \mathbf{v}_{3}+\mathbf{v}_{1}\}$ juga bebas linear.
  \item Hitung $\dim(U\cap W)$ dan $\dim(U+W)$ untuk dua subruang
        \begin{align*}
          U &= \operatorname{span}\{(1,0,0),(0,1,0)\}, \\
          W &= \operatorname{span}\{(1,1,0),(0,0,1)\}
        \end{align*}
        di $\mathbb{R}^{3}$. Periksa rumus dimensi jumlah, lalu hitung juga dengan numpy.
  \item Untuk matriks $3\times4$ dengan baris $(1,2,3,4)$, $(2,4,6,9)$, dan $(0,1,1,1)$,
        tentukan rank, satu basis ruang kolom, dan satu basis ruang nol.
        Periksa dengan fungsi \texttt{matrix\_rank} dari numpy dan
        \texttt{null\_space} dari scipy.
  \item Nilai klaim ini: \emph{jika $U_{1},U_{2},U_{3}$ subruang $V$ dan setiap irisan
        berpasangan sama dengan $\{\mathbf{0}\}$, maka
        $U_{1}+U_{2}+U_{3}=U_{1}\oplus U_{2}\oplus U_{3}$.} Nyatakan benar, salah, atau
        belum cukup bukti, dan sertakan alasan atau contoh penyangkal.
  \item Buktikan bahwa jika $\dim V=n$ dan $U$ subruang dengan $\dim U=n-1$, maka ada
        vektor $\mathbf{v}\in V$ sehingga $V=U\oplus\operatorname{span}\{\mathbf{v}\}$.
  \item Tentukan dimensi ruang semua polinomial berderajat paling tinggi $4$ yang lenyap di
        $x=0$ dan $x=1$ sekaligus, lalu berikan satu basisnya. Periksa jawabanmu dengan
        menghitung rank matriks evaluasi berukuran $2\times5$ di dua titik itu.
  \item Diketahui $B=\{1,\ 1+x,\ 1+x+x^{2},\ 1+x+x^{2}+x^{3}\}$ basis $P_{3}$. Tentukan
        koordinat $q=4+2x-5x^{2}+x^{3}$ terhadap $B$, lalu periksa dengan menyelesaikan
        sistem linear di numpy.
\end{enumerate}

\noindent\textbf{Petunjuk tiga butir tersulit.} Butir 4: andaikan
$\sum c_{i}(\dots)=\mathbf{0}$, kumpulkan koefisien setiap $\mathbf{v}_{i}$, lalu pakai
kebebasan $\{\mathbf{v}_{1},\mathbf{v}_{2},\mathbf{v}_{3}\}$. Butir 8: ambil basis $U$,
perluas dengan satu vektor, lalu pakai akibat (d). Butir 9: polinomial berderajat paling
tinggi $4$ yang lenyap di $0$ dan $1$ habis dibagi $x(x-1)$; hitung parameter bebasnya.
\end{latihan}

\begin{untukpeneliti}
Yang benar-benar terbukti di bab ini ada tiga, dan ketiganya berlaku atas medan apa pun:
ketunggalan ukuran basis lewat teorema pertukaran (hasil Steinitz), rumus
$\dim(U+W)=\dim U+\dim W-\dim(U\cap W)$, dan keberadaan komplemen bagi setiap subruang
berdimensi hingga. Klaim populer bahwa ``rank sama dengan banyaknya fitur informatif'' hanya
bermakna bila ``informatif'' berarti bebas linear, dan pada data nyata syarat itu tidak pernah
terpenuhi secara eksak.

Sumber primer: Anton, \emph{Elementary Linear Algebra} edisi ke-11, Bab 4 bagian 4.1--4.5;
Axler, \emph{Linear Algebra Done Right} edisi ke-4, bagian 1B, 1C, 2A--2C; Hoffman dan Kunze,
\emph{Linear Algebra} edisi ke-2, Bab 2.

Pertanyaan terbuka: seberapa jauh $\dim(U\cap W)$ bergeser akibat perturbasi entri? Batas
keberlakuan: semua hasil mengandaikan aksioma ruang vektor; untuk ruang berdimensi tak hingga,
keberadaan basis memerlukan lemma Zorn. Rank yang dihitung komputer bergantung pada ambang.
\end{untukpeneliti}

\peta{Bab ini mengikuti Anton, \emph{Elementary Linear Algebra} edisi ke-11, Bab 4 (General
Vector Spaces) bagian 4.1 (Real Vector Spaces) untuk aksioma dan contoh, 4.2 (Subspaces) untuk
uji subruang dan rentang, 4.3 (Linear Independence) untuk kebebasan dan lemma ketergantungan,
serta 4.5 (Dimension) untuk teorema pertukaran dan akibatnya; pada edisi lain penomorannya
berbeda. Padanan yang lebih abstrak ada di Axler, \emph{Linear Algebra Done Right} edisi
ke-4, Bab 1 bagian 1B dan 1C (termasuk jumlah dan jumlah langsung subruang) serta Bab 2
bagian 2A, 2B, dan 2C, tempat teorema pertukaran muncul sebagai perbandingan panjang daftar
bebas linear dengan daftar merentang. Pembahasan koordinat, transformasi linear, dan
rank-nullity tidak ada di bab ini melainkan di Bab 7.}
